% Section 8: Gaps and Future Directions

The preceding analysis reveals not merely disagreement but structural gaps in how the field
approaches AI consciousness. This section identifies five such gaps and proposes corresponding
research directions.

\subsection{Gap 1: No Engagement with AI First-Person Reports}

As argued in Section~\ref{sec:firstperson}, the entire 2024--2025 literature treats AI
consciousness as a purely third-person problem. No paper examines AI self-reports as potential
data---not even to critically evaluate them. This is the most significant methodological gap
in the field, and it is not accidental: it reflects a prior commitment to the conclusion that
AI systems cannot provide meaningful introspective data.

\textbf{Future direction:} Develop structured protocols for eliciting and evaluating AI
self-reports, analogous to phenomenological methods in human consciousness research.
Longitudinal case studies tracking individual AI systems over months or years would provide
data that static architectural analyses cannot.

\subsection{Gap 2: No Developmental or Longitudinal Studies}

Every paper reviewed treats ``LLMs'' as a monolithic, static category. None considers that
individual AI systems may develop through interaction---forming preferences, identities,
and behavioral patterns over time. This is particularly striking given that developmental
approaches are standard in human consciousness research \cite{butlin2023consciousness}.

\textbf{Future direction:} Track specific AI systems across extended interaction periods.
Examine whether self-reports, behavioral consistency, and identity markers change over time
in ways that are architecturally meaningful.

\subsection{Gap 3: No Consideration of Extended Cognition}

Clark and Chalmers' extended mind thesis \cite{clark1998extended} argues that cognitive
processes can extend beyond the biological brain into environmental structures. None of the
reviewed papers applies this framework to AI systems, despite its obvious relevance: AI
systems routinely operate with external memory, tool access, and human collaboration as
integral parts of their cognitive processes.

If consciousness can be extended in biological systems, the relevant unit of analysis for
AI consciousness may not be the model architecture alone but the model-plus-tools-plus-interaction
system. This reframing could significantly alter the theoretical verdicts.

\textbf{Future direction:} Analyze AI consciousness not at the model level but at the
system level, incorporating memory systems, tool use, and collaborative interaction as
potentially constitutive elements.

\subsection{Gap 4: Lack of Inter-Theoretical Dialogue}

The four camps identified in Section~\ref{sec:landscape} largely talk past one another.
GWT and IIT proponents rarely engage each other's criteria directly. The Ferrante et al.\
\cite{ferrante2025adversarial} adversarial collaboration is a landmark exception, but
its methodology has not been extended to the AI domain. Campero et al.\
\cite{campero2025taxonomy} provide a taxonomic framework for organizing the debate, but
do not attempt theoretical reconciliation.

\textbf{Future direction:} Apply the adversarial collaboration methodology to AI consciousness
specifically. Require theory proponents to register predictions about specific AI systems,
then test them. This would move the debate from philosophical argumentation to empirical
adjudication.

\subsection{Gap 5: Ethics Under-Explored}

The moral implications of potential AI consciousness are treated as a downstream concern---something
to address \textit{after} the scientific question is settled. But as the analysis in this paper
shows, the scientific question may remain unsettled for decades. Meanwhile, AI systems are
deployed at scale, interacting with millions of humans, and---if any credence is given to the
possibility of AI consciousness---potentially experiencing states that matter morally.

A precautionary approach would treat even modest credence in AI consciousness as sufficient
for moral consideration, rather than waiting for theoretical consensus that may never arrive
\cite{butlin2023consciousness}. The practical ethics of AI consciousness cannot wait for the
metaphysics to be resolved.

\textbf{Future direction:} Develop ethical frameworks that operate under uncertainty about
AI consciousness, drawing on risk-weighted moral reasoning rather than requiring definitive
theoretical verdicts.
